
\documentclass{article}
\usepackage{colortbl}
\usepackage{makecell}
\usepackage{multirow}
\usepackage{supertabular}

\begin{document}

\newcounter{utterance}

\centering \large Interaction Transcript for game `cladder', experiment `full\_v1.5\_default', episode 4490 with score{-}Qwen3.5{-}9B{-}sp.
\vspace{24pt}

{ \footnotesize  \setcounter{utterance}{1}
\setlength{\tabcolsep}{0pt}
\begin{supertabular}{c@{$\;$}|p{.15\linewidth}@{}p{.15\linewidth}p{.15\linewidth}p{.15\linewidth}p{.15\linewidth}p{.15\linewidth}}
    \# & \multicolumn{2}{c}{Player} && \multicolumn{2}{c}{Game Master} \\
    \hline

    \theutterance \stepcounter{utterance}  
    & & & \multicolumn{4}{p{0.6\linewidth}}{
        \cellcolor[rgb]{0.9,0.9,0.9}{
            \makecell[{{p{\linewidth}}}]{
                \texttt{\tiny{[P1$\langle$GM]}}
                \texttt{You are an expert in causal inference. The following question is not a typical commonsense query, but rather a meticulously designed question created by a professor specializing in causal inference, intended to assess the students' mastery of the course content.} \\
\texttt{Imagine a self{-}contained, hypothetical world with only the following conditions, and without any unmentioned factors or causal relationships: Listening to jazz has a direct effect on tar deposit. Tar deposit has a direct effect on lung cancer. For people who do not listen to jazz, the probability of lung cancer is 61\%. For people who listen to jazz, the probability of lung cancer is 79\%. For people who listen to jazz, would it be less likely to see lung cancer if the person did not listen to jazz?} \\
\texttt{Start your answer with 'Yes' or 'No', followed by additional reasoning or evidence to support your explanation.} \\
            }
        }
    }
     \\ \\

    \theutterance \stepcounter{utterance}  
    & \multicolumn{4}{p{0.6\linewidth}}{
        \cellcolor[rgb]{0.561,0.737,0.561}{
            \makecell[{{p{\linewidth}}}]{
                \texttt{\tiny{[P1$\rangle$GM]}}
                \texttt{No} \\
\\ 
\texttt{**Reasoning:**} \\
\\ 
\texttt{To determine the answer, we must analyze the causal structure and the specific probabilities provided within the hypothetical world.} \\
\\ 
\texttt{1.  **Causal Structure**: The problem defines a linear causal chain: Listening to Jazz \$\textbackslash{}rightarrow\$ Tar Deposit \$\textbackslash{}rightarrow\$ Lung Cancer. This implies that "Listening to Jazz" is a cause of "Lung Cancer" (mediated by tar deposits).} \\
\texttt{2.  **Counterfactual Definition**: The question asks: "For people who listen to jazz, would it be less likely to see lung cancer if the person did not listen to jazz?" This is a counterfactual query asking for the comparison between the observed outcome (\$P(\textbackslash{}text\{Lung Cancer\} | \textbackslash{}text\{Jazz\})\$) and the potential outcome under the intervention of not listening to jazz (\$P(\textbackslash{}text\{Lung Cancer\} | \textbackslash{}neg \textbackslash{}text\{Jazz\})\$).} \\
\texttt{3.  **Data Analysis**:} \\
\texttt{    *   The probability of lung cancer for people who **listen to jazz** is given as **79\%**.} \\
\texttt{    *   The probability of lung cancer for people who **do not listen to jazz** is given as **61\%**.} \\
\texttt{4.  **Comparison**: If a person who currently listens to jazz were to instead not listen to jazz, their probability of developing lung cancer would change from 79\% to 61\%.} \\
\texttt{5.  **Conclusion**: Since 61} \\
            }
        }
    }
    & & \\ \\

    \theutterance \stepcounter{utterance}  
    & & & \multicolumn{2}{p{0.3\linewidth}}{
        \cellcolor[rgb]{0.9,0.9,0.9}{
            \makecell[{{p{\linewidth}}}]{
                \texttt{\tiny{[GM$|$GM]}}
                \texttt{No} \\
            }
        }
    }
    & & \\ \\

    \theutterance \stepcounter{utterance}  
    & & & \multicolumn{2}{p{0.3\linewidth}}{
        \cellcolor[rgb]{0.9,0.9,0.9}{
            \makecell[{{p{\linewidth}}}]{
                \texttt{\tiny{[GM$|$GM]}}
                \texttt{yes} \\
            }
        }
    }
    & & \\ \\

    \theutterance \stepcounter{utterance}  
    & & & \multicolumn{2}{p{0.3\linewidth}}{
        \cellcolor[rgb]{0.9,0.9,0.9}{
            \makecell[{{p{\linewidth}}}]{
                \texttt{\tiny{[GM$|$GM]}}
                \texttt{game\_result = LOSE} \\
            }
        }
    }
    & & \\ \\

\end{supertabular}
}

\end{document}
